\section{How We Evaluate}
\label{sec:method}

A volunteer platform with thousands of machines, controlled churn, hostile owners and paired comparisons of policies cannot be run on real hardware. We evaluate \voss{} in a deterministic simulator and calibrate it against measurements. This section says what the simulator models, where its numbers come from, how far we checked it, and how every comparison is made.

\subsection{The simulator}
\label{sec:sim}
The simulator is a single-threaded discrete-event simulator in Python. Simulated time jumps from one event to the next in a fixed order, and every random draw comes from a named, seeded stream, so a run can be repeated exactly.
A run takes a configuration, a calibration bundle and a seed. It first draws everything the system under test cannot influence: the machines, every session start and end, which owners are adversaries, the arrivals and the noise in each execution. These traces are content-hashed, so two systems compared on one seed face the same world.
The engine then runs the actors (workers, schedulers, routers, orchestrator, verifier, key services) and the policies. Policies see only immutable views built from messages, never the world itself. At the end, checks assert that every success ran on an attested host after a key release, that each invocation has one outcome, and that the time components add up to the measured latency. Analysis reads only the raw tables.

\textbf{Nothing is free.} Every interaction is a message that arrives after a delay from the distance between about 100 cities (100\,km per millisecond, times 1.3 for route detours), plus last-mile delay and jitter. Disk reads, CPU contention, TPM quotes, key release and the decision time of planner, placer, router and orchestrator all consume simulated time.
Decision time is a fixed part plus a part proportional to the work, so results do not depend on the speed of the machine running the simulation. The coefficients are placeholders from the literature, not measurements. For the planner the charge is about 0.1\,s, for the greedy and the exact planner alike, while the exact planner really takes \MilpSolveS\,s, so its results are optimistic about how stale its plans are.

\textbf{The world.} Workers come in three tiers (nightly, idle-hours, dedicated) with different hardware, cities and trust substrates. Session lengths are heavy-tailed with time-of-day effects, most departures are abrupt, and correlated shocks remove a share of a city at once.
No public trace of consumer donation sessions exists, so we follow volunteer-availability studies~\cite{javadi2013fta}, mark the parameters \emph{assumed} and sweep them. The workload has six tenants and 12 functions modelled on SeBS~\cite{copik2021sebs}, in three service classes, with skewed popularity and bursty, day-and-night arrivals~\cite{shahrad2020wild}. 11 of the 12 functions are marked idempotent, by our assumption.
A fraction $\alpha$ of workers are adversaries who omit, delay, black-hole, defect after building a record, send fake drain notices, or collude within a failure domain. By Property 1 none of them forges results or reads keys on an attested substrate, while the unattested reference allows both.

\subsection{Calibration and what we checked}
\label{sec:calib}
Parameters from the specification are literature values or assumptions and are tagged so. Measured values come only from calibration bundles. We measured one on a consumer laptop (\CalCpu, a discrete TPM~2.0, a hard disk), which stands in for a donor machine.
A harness reboots it into the appliance and records timestamps at about one-second resolution, and microbenchmarks cover TPM quotes, hashing, microVM boot and snapshots. Compared with our placeholders, reboot to kernel is slower (\CalReboot\,s against \CalRebootPh\,s), the phase from kernel start to a running agent is far slower (\CalUserspace\,s), and creating a 512\,MB snapshot on a hard disk takes \CalSnapCreate\,s.
From power-on, an attested worker takes \CalAttested\,s (median of \CalAttestedN{} boots, \CalAttestedKernel\,s of it after kernel start). A warm-cache microVM boot takes \CalFcBootWarmMs\,ms, close to published figures~\cite{agache2020firecracker}.

Two limits apply. The laptop is not the intended donor machine and lacks AMD SEV-SNP and Intel TDX, so confidential-VM parameters are literature values, which we sweep~\cite{misono2025cvm,holmes2024severifast,li2023bifrost}.
And the appliance ran in \emph{host mode} (the measurement chain and quote on the host, microVMs through Firecracker) rather than as a locked image on bare metal. The attested-session figure is a measured sum, not a measured attestation of the final image.

\textbf{A first check against real execution (E01).} To test the start-up and execution parts of the simulator, we ran the same scenarios in a single-box emulation on a lab node (Intel Core i7-12700, 125\,GB, NVMe, Firecracker 1.10.1) and in the simulator, with a calibration measured on that node. The emulation uses real microVMs for boot, snapshot restore, dispatch and function execution, and keeps network, key release, attestation and queueing as models (16 donor slices, five functions, ten minutes per run, three seeds).
Success rate agrees within 0.6 percentage points in every scenario, the mix of start-up levels is identical, and median synchronous latency agrees within 6--23\%. The \emph{shape} of the latency distribution does not meet our pre-set cut-offs, so every scenario is graded ``partly'' close and none ``close''. At light load the simulator's P99 is 3.7 to 4.1 times the emulation's, driven by a handful of cold invocations: the P99 of about 530 requests is the five slowest.
E01 does not test the network, attestation, key release, session boot or failure detection at scale, and no multi-machine deployment has validated the simulator end to end. We therefore read differences between designs, not absolute latencies.

\subsection{Protocol}
\label{sec:protocol}
Each experiment is a frozen specification, and every cell is run with ten seeds (three for the two largest scalability points). Unless stated, \voss{} uses the greedy planner and the default placer, router and hedger. Systems in a comparison share the world and the per-invocation random draws, so each comparison is \emph{paired} by seed.
For a metric we take the per-seed difference and report its mean with a 95\% bootstrap interval~\cite{efron1993bootstrap}. Percentiles are computed per run, then averaged, never pooled.

Each hypothesis is written down before the run with a mechanical rule. It is \emph{supported} if the interval excludes zero in the predicted direction by at least a stated margin, \emph{weak} if it excludes zero but the effect is below the margin, \emph{inconclusive} if the interval includes zero, and \emph{refuted} if it excludes zero the other way.
The manuscript's numbers and tables are generated from the result tables by a script that also fails if a result stops supporting the sentence it backs.

We report success, the share of deadlines met, median and P99 synchronous latency, \emph{waste} (extra attempts per success), donor memory held warm, and key exposure. A system is \emph{feasible} in a scenario cell if the lower 95\% bound of its success is at least 0.99 and either its deadline-met share or its P99 meets a target.
Feasibility maps then show \emph{where} each design works, rather than an average.

\textbf{Baselines} run on the same simulator, candidates and cost model and differ only in the registered policy: no planning (provision on arrival), keep-alive with a time limit, a heartbeat-only detector, and $f=0$ (no robustness). The \emph{simple baseline} combines static regional clusters, nearest-region routing, least-loaded placement, keep-alive and no hedging.
Commercial-style references (\emph{Lambda-like}, \emph{Azure-like}) are always-on datacenter hosts with per-function keep-alive and nearest-region routing, priced from list prices, on the same engine.

\subsection{Reproducibility}
\label{sec:repro}
Every number, table and figure in Section~\ref{sec:results} is regenerated from raw run logs by committed scripts, and runs are bit-reproducible: re-executing a run gives raw tables with the same result hash, recorded in a manifest with the scenario and calibration hashes. Floating-point arithmetic repeats bit for bit only under stated conditions, because a last-bit difference can flip a near-tie in placement and send the whole trajectory elsewhere.
We found that NumPy's maths functions use different code paths on CPUs with and without AVX-512, and a first verification re-run failed for exactly this reason. We now pin NumPy to the SSE4 baseline that every x86-64 CPU supports and record a canary hash of the maths functions in each manifest.
The full sweep, 21 experiments and 2{,}366 runs, was then re-executed on independently provisioned machines, and 2{,}366 of 2{,}366 hashes matched.
Reproduction needs an x86-64 CPU with fused multiply-add, Linux with glibc, Python 3.11 and the committed lockfile, and about four hours of wall-clock time as parallel jobs. No TPM or special hardware is needed. Appendix~\ref{app:repro} gives the requirements, checks and commands.
